% !TEX root = ../main.tex
\section{Scorecard against the proposal}
\label{sec:scorecard}

Table~\ref{tab:scorecard} sets each objective of the proposal against the result of this build. The status words are decided from the outputs present when the specification is generated.

\begin{table}[htbp]
\centering
\caption{Proposal objectives and the result in this build.}
\label{tab:scorecard}
\footnotesize
\begin{tabularx}{\linewidth}{@{}>{\raggedright\arraybackslash}p{0.27\linewidth} L l >{\raggedright\arraybackslash}p{0.24\linewidth}@{}}
\toprule
objective (proposal) & result in this build & status & note \\
\midrule
Province-week panel 2016--2020 on PSGC: dengue, rain, Sentinel-1, hazard, storms & \nUnits{} harmonised units $\times$ \nWeeks{} weeks; every listed layer plus river routing & \statusPanel & \nUnits{} harmonised 2016--2020 units, not \proposalProvinces{} provinces (\dgUnits{} reporting units); Sentinel-1 only in the four storm windows \\
Lead-lag graph (PCMCI+), sentinels (PageRank), communities (Louvain); test sentinel features out of sample & \llEdges{} links at $q < \fdrQ$; sentinels and communities on the dashboard; ablation run & \stOk & tested and null: graph against no graph, relative WIS \relWISGraph{} \\
Relative WIS below 1 against seasonal naive, 2018--2019 & XGBoost \relWIS{}; median ensemble \ensRelWIS{} & \stOk & \\
90\% interval coverage of \targetCov{} & \covNinety{} overall; \covNinetyRange{} by horizon & \stOk & four-week horizon close to the floor \\
Benchmarks: Chronos-2, TabPFN-TS, MPNN-LSTM; median ensemble; conformal check & all three and the ensemble scored on the same rows; online conformal check & \statusBench & the ensemble is evaluated, but storm scenarios use XGBoost (the zero-shot models take no storm-track input) \\
Access loss per municipality, friction-surface and network methods, four typhoons & both methods, \nAdmThree{} municipalities, four storms $\times$ \nTracks{} tracks & \statusAccess & method disagreement reported as the $A_{\text{lo}}$--$A_{\text{hi}}$ range \\
Sensitivity bounds (Kendall's $\tau$); Metro Manila benchmark & two methods, $\tau \geq \tauMinAll$; NCR \benchNCR{} against \litNCRShare{} & \stOk & order-of-magnitude check \\
Dashboard: ranked provinces, calibrated risk, access, sentinels, drivers, tier, any historical or shifted track & \nScen{} backtest scenarios and \nCat{} catalogue storms, pre-computed; $R$ recalibrated & \stOk & tiers are provisional rules \\
Backtest: did high-priority provinces see more cases? & $\rho(R)$ \rhoRRange{} over four storms & \stOk & 2020 storms indicative only (COVID-19 reporting) \\
System Usability Scale from three target users & test form and task built into the dashboard; \susN{} responses so far & \statusSUS & needs sessions with health and disaster-office users \\
Leptospirosis (stretch): national lag analysis; region-week series if digitised & national lag analysis, an open \leptoNCities{}-city series and a typhoon event study & \statusLepto & association only, not in the forecast; no province-level series \\
\bottomrule
\end{tabularx}
\end{table}

\subsection{Proposal corrections found in the data}

\begin{itemize}
\item \dgUnits{} dengue reporting units, not \proposalProvinces{} provinces (harmonised to \nUnits{}); deaths are unreliable; 2020 collapsed under COVID-19.
\item An open subnational leptospirosis series exists (Project CCHAIN, \leptoNCities{} cities).
\item The EMS activation for Rolly has damage grading only, no flood delineation; the Global Flood Database has no Philippine event after 2017 \parencite{tellman2021}.
\item The flood hazard maps come from the NOAH archive at 5 and 25 years, not the national lidar portal; \nHazNoMap{} units have no 25-year map.
\item HydroRIVERS places the Philippines in the Australasia file, not the Asia file.
\item The ferry routes describe 2024--2026 operations with partly coarse port coordinates; the official facility registry blocks scripted download (the OSM and NOAH fallback is used).
\item HRSL under-counts people aged 60 or over against the 2020 census.
\item The proposal's final forecast was the median ensemble; storm scenarios use XGBoost, the only storm-conditioned model.
\end{itemize}
